% ============================================================
% UNIDAD 7: OPERACIONES AÉREAS, ESPACIO AÉREO Y CARTAS
% ============================================================
\clearpage
\section*{UNIDAD 7: OPERACIONES AÉREAS}
\subsection*{Tipos de operación, espacio aéreo y cartas aeronáuticas}

\begin{recuadro}
\textbf{Objetivo de la unidad:} conocer cómo se organizan las operaciones aéreas, cómo se divide y
controla el espacio aéreo y cómo se leen las cartas aeronáuticas, para entender el entorno en el
que funcionan los equipos de comunicaciones, navegación y vigilancia que mantiene el técnico
aviónico.
\end{recuadro}

% ------------------------------------------------------------
\section*{1. TIPOS DE OPERACIÓN}

\begin{center}
\small
\begin{tabular}{|p{0.30\linewidth}|p{0.14\linewidth}|p{0.44\linewidth}|}
\hline
\textbf{Tipo} & \textbf{Norma} & \textbf{Ejemplos} \\ \hline
Aviación general & RAAC 91 & Vuelos privados, aeroclubes, instrucción, aviación corporativa. \\
\hline
Transporte aéreo regular & RAAC 121 & Líneas aéreas con aviones grandes. \\ \hline
Transporte no regular y operaciones menores & RAAC 135 & Taxis aéreos, ambulancias aéreas,
vuelos ejecutivos. \\ \hline
Trabajo aéreo & RAAC 137 & Aviación agrícola, fotografía, inspección de líneas eléctricas. \\
\hline
Aeronaves pilotadas a distancia & RAAC 100 & Drones (RPAS). \\ \hline
Aviación de Estado & -- & Militar, policial, sanitaria oficial (matrícula LQ). \\ \hline
\end{tabular}
\end{center}

\paragraph*{Reglas de vuelo.}
\begin{itemize}
  \item \textbf{VFR} (reglas de vuelo visual): el piloto se separa de otras aeronaves y del
        terreno mirando hacia afuera. Requiere condiciones meteorológicas visuales (VMC).
  \item \textbf{IFR} (reglas de vuelo por instrumentos): se vuela con referencia a los
        instrumentos y siguiendo procedimientos y autorizaciones del control de tránsito. Es la
        regla de la aviación comercial, y la que exige que la aviónica de navegación y
        comunicaciones funcione con total confiabilidad.
\end{itemize}

% ------------------------------------------------------------
\section*{2. LOS SERVICIOS DE TRÁNSITO AÉREO}

En la Argentina los presta la \textbf{EANA}. Incluyen:
\begin{itemize}
  \item \textbf{Control de tránsito aéreo (ATC):} \textbf{torre de control} (TWR) para el aeródromo
        y sus inmediaciones; \textbf{control de aproximación} (APP) para las llegadas y salidas; y
        \textbf{centro de control de área} (ACC) para la fase en ruta.
  \item \textbf{Servicio de información de vuelo} (FIS): datos útiles para la seguridad del vuelo.
  \item \textbf{Servicio de alerta}: notifica a búsqueda y salvamento cuando una aeronave necesita
        ayuda.
\end{itemize}
Para prestarlos, el espacio aéreo argentino se divide en \textbf{regiones de información de vuelo}
(FIR), cada una con su centro de control de área.

\paragraph*{El transpondedor.} Es el equipo de a bordo que responde a los radares secundarios
con un código de cuatro dígitos asignado por el control. Algunos códigos son universales:
\textbf{7500} (interferencia ilícita), \textbf{7600} (falla de comunicaciones) y \textbf{7700}
(emergencia). Con \textbf{ADS-B}, la aeronave además transmite su posición GPS, su altitud y su
velocidad.

% ------------------------------------------------------------
\section*{3. EL ESPACIO AÉREO}

\subsubsection*{3.1. Espacios aéreos controlados y no controlados}
\begin{itemize}
  \item \textbf{Zona de control (CTR):} rodea a un aeródromo controlado, desde el suelo hasta un
        límite superior.
  \item \textbf{Área de control terminal (TMA):} encima y alrededor de una o varias CTR, para
        ordenar llegadas y salidas.
  \item \textbf{Aerovías (AWY):} ``rutas'' del cielo, con ancho y altitudes definidos, que unen
        radioayudas o puntos de notificación.
  \item \textbf{Espacios reglamentados:} \textbf{prohibidos} (no se puede volar),
        \textbf{restringidos} (solo con condiciones) y \textbf{peligrosos} (hay actividades de
        riesgo, como tiro o vuelo de práctica militar).
\end{itemize}

\subsubsection*{3.2. Clases de espacio aéreo (OACI)}
\begin{center}
\small
\begin{tabular}{|c|p{0.18\linewidth}|p{0.64\linewidth}|}
\hline
\textbf{Clase} & \textbf{Vuelos permitidos} & \textbf{Servicio} \\ \hline
A & Solo IFR & Todos controlados y separados entre sí. \\ \hline
B & IFR y VFR & Todos controlados y separados entre sí. \\ \hline
C & IFR y VFR & IFR separados de todos; VFR separados de IFR y con información de otros VFR. \\
\hline
D & IFR y VFR & IFR separados de IFR; todos reciben información de tránsito. \\ \hline
E & IFR y VFR & IFR controlados y separados de IFR; VFR no controlados. \\ \hline
F & IFR y VFR & Asesoramiento de tránsito para IFR (no controlado). \\ \hline
G & IFR y VFR & No controlado; solo información de vuelo. \\ \hline
\end{tabular}
\end{center}

% ------------------------------------------------------------
\section*{4. LAS CARTAS AERONÁUTICAS}

Las cartas aeronáuticas se rigen por el \textbf{Anexo 4} de la OACI y se publican, junto con toda la
información oficial para volar en el país, en la \textbf{AIP} (Publicación de Información
Aeronáutica) de la Argentina. Los cambios temporarios se informan por \textbf{NOTAM} (aviso a los
aviadores): por ejemplo, una radioayuda fuera de servicio por mantenimiento.

\begin{center}
\small
\begin{tabular}{|p{0.32\linewidth}|p{0.58\linewidth}|}
\hline
\textbf{Carta} & \textbf{Para qué sirve} \\ \hline
Carta aeronáutica visual (escala 1:500.000 o 1:1.000.000) & Vuelo VFR: relieve, ciudades, rutas,
ríos, aeródromos, espacios aéreos. \\ \hline
Carta de navegación en ruta (ENR) & Vuelo IFR: aerovías, radioayudas, altitudes mínimas, límites
de FIR. \\ \hline
Salidas y llegadas normalizadas (SID y STAR) & Procedimientos para salir de un aeródromo y para
llegar a él. \\ \hline
Carta de aproximación por instrumentos (IAC) & Procedimiento para aproximarse y aterrizar con
ILS, VOR, RNAV (GPS). \\ \hline
Plano de aeródromo (ADC) & Pistas, calles de rodaje, plataformas, frecuencias. \\ \hline
\end{tabular}
\end{center}

\subsubsection*{4.1. Qué se lee en una carta}
\begin{itemize}
  \item \textbf{Coordenadas} (latitud y longitud), escala y proyección.
  \item \textbf{Rumbos} magnéticos y la \textbf{declinación magnética} (diferencia entre el norte
        verdadero y el magnético).
  \item \textbf{Radioayudas}, con su símbolo, nombre, identificación en código Morse y
        \textbf{frecuencia}: \textbf{VOR} (VHF, 108 a 117,95 MHz), \textbf{DME} (distancia),
        \textbf{NDB} (baja/media frecuencia), \textbf{ILS} (localizador y senda de planeo).
  \item \textbf{Altitudes mínimas} de las rutas y del terreno.
  \item \textbf{Frecuencias de comunicaciones} de torre, aproximación y control.
  \item \textbf{Espacios aéreos} con sus clases y límites verticales.
\end{itemize}

\begin{recuadro}
\textbf{¿Por qué le importa al técnico aviónico?} Cada símbolo de la carta es un sistema que el
técnico mantiene o del que depende un equipo de a bordo: el receptor VOR/ILS, el DME, el ADF, el
GPS, el transpondedor y las radios de VHF. Entender cómo los usa el piloto ayuda a diagnosticar
una falla reportada (``el VOR 1 marca 5 grados de diferencia con el VOR 2'') y a dimensionar su
impacto en la operación y en la lista de equipo mínimo (MEL).
\end{recuadro}

% ------------------------------------------------------------
\section*{5. ACTIVIDADES DE LA UNIDAD}
\begin{enumerate}
  \item \textbf{Lectura de cartas:} con una carta visual y una de aproximación de un aeródromo
        cercano (por ejemplo, Aeroparque, Ezeiza o San Fernando), identificar radioayudas,
        frecuencias, espacios aéreos y altitudes mínimas.
  \item \textbf{NOTAM:} buscar NOTAM vigentes de un aeródromo y explicar cuáles afectan equipos de
        navegación o comunicaciones.
  \item \textbf{De la carta al equipo:} para una aproximación ILS, listar todos los equipos de a
        bordo y de tierra que intervienen y qué pasaría si fallara cada uno.
  \item \textbf{Clases de espacio:} para tres vuelos (un VFR de aeroclub, un vuelo de línea, un
        drone), indicar en qué espacios pueden volar y qué necesitan.
\end{enumerate}
